%! Composition of Functions — Unit 1, Lesson 1.5 — Warm-Up
\documentclass[10pt]{article}
\usepackage{precalculus-article}
\usepackage{precalculus-boxes}
\newcommand{\TallMath}[1]{$\displaystyle #1\rule[-1.4em]{0pt}{3.2em}$}

\begin{document}

\pageheader{Unit 1, Lesson 1.5}{Warm-Up}

\vspace{0.12in}

\begin{enumerate}[itemsep=10pt]

\item From Lesson 1.1: let $f(x) = 2x + 1$. Whatever goes into the parentheses
      is what gets doubled --- even if it is not a number.

\begin{work}
  f(4) &= 2(4) + 1 \\
       &= 9 \\
  f(-3) &= 2(-3) + 1 \\
        &= -5 \\
  f(x^2) &= 2(x^2) + 1 \\
         &= 2x^2 + 1
\end{work}

\begin{enumerate}[label=(\alph*), itemsep=4pt]
  \item $f(4) = $ \blank{2.4cm} \qquad $f(-3) = $ \blank{2.4cm}
  \item $f(x^2) = $ \blank{3.4cm} \quad (small question today, big payoff later)
\end{enumerate}

\item From Lesson 1.1: the table gives every value of a function $g$.

\smallskip
\renewcommand{\arraystretch}{1.4}
\begin{tabular}{|c|c|c|c|c|c|}
\hline
$x$ & 0 & 1 & 2 & 3 & 4 \\
\hline
$g(x)$ & 3 & 0 & 4 & 1 & 2 \\
\hline
\end{tabular}
\smallskip

\begin{enumerate}[label=(\alph*), itemsep=4pt]
  \item $g(2) = $ \blank{2cm} \qquad $g(4) = $ \blank{2cm}
  \item For which input $x$ is $g(x) = 3$? \quad $x = $ \blank{2cm}
\end{enumerate}

\item From Lesson 1.4: in the rule $g(x) = f(x + 3)$, the ``$+3$'' is written
      \textbf{outside} / \textbf{inside} (circle one) the function $f$, so it
      changes the \blank{2.4cm} rather than the outputs.

\end{enumerate}

\end{document}
